% Ad Quintum Fratrem 3.3 --- headnote
% ad-quintum-fratrem-03-03, 1 November 54 BC, Rome

\textit{Cicero to Quintus, written from Rome on or about the
Kalends of November 54 BC. The dateline in the manuscripts is
corrupt (\textit{a. d. xit K. Nov.}); modern editors read the
Kalends themselves --- 1 November --- on the body of the
letter's evidence, with the Gabinius \textit{maiestas} verdict
still expected ``within three days'' (\textsection 3). Four
sections, dictated to a secretary at the end of a day spent
defending in court.}

\textit{The first section is the private side of the letter
--- the painful interval, more than fifty days, in which not
a letter and not even a rumour from Quintus, from Caesar, or
from northern Gaul had reached Rome. Caesar's second British
campaign was being concluded in October; the Channel crossing
back to Gaul and the winter dispersal of the legions were
imminent. Cicero's anxiety about ``the sea over there and the
land'' is the running undercurrent of the year's letters; the
news that came up to Rome in the next weeks --- Quintus and
his legion of recruits would later be besieged in his winter
camp by Ambiorix and the Eburones --- would justify the
worry.}

\textit{The political news in \textsection 2 is the famous
four-candidate ambitus hurricane that the previous letter
(Q. fr. 3.2) had inaugurated. By 1 November every man
standing for the consulship of 53 had been prosecuted for
canvassing; the elections were being knocked off, one day at
a time, through formal announcements of bad omens
(\textit{obnuntiationes}), with the entire \textit{boni}
class delighted, because the consuls of 54 were under
universal suspicion of having taken bribes from the
candidates to fix the result. The pact (exposed by Memmius
in the Senate on 30 September) had been a four-way collusion
in which Memmius and Calvinus, the candidates favoured by
Caesar's money, agreed with the consuls Appius Claudius and
Domitius Calvinus to deliver Memmius and Calvinus to the
fasces in exchange for fraudulent senatorial votes after the
fact. Of the four candidates Cicero is now most concerned
with M. Valerius Messala Rufus, whom he would later defend
on the ambitus charge; ``our Messala safe'' is, as he says,
``linked also to the safety of the rest.'' Gabinius
prosecuted on top of his \textit{maiestas} ordeal: P. Sulla
the former candidate of 66 BC, his stepson C. Memmius
(distinct from the consular candidate), his brother Caecilius,
and his son Faustus Sulla flanking the prosecution. L.
Torquatus, who argued against (probably defending the
substantive prosecution), failed despite popular goodwill.}

\textit{\textsection 3 reports on the impending \textit{maiestas}
verdict; this letter precedes Q. fr. 3.4 by a few days, and
when the news of acquittal (32 to 70) breaks Cicero will
write a fuller post-mortem in the next letter. Cn. Domitius
Calvinus is the consular candidate of the year sitting as
juror in the \textit{maiestas} trial; Alfius the president of
the court is the man whose edict on \textit{maiestas} brought
Gabinius back to Rome in Q. fr. 3.1 \textsection 24.
\textsection 4 closes the letter on the boy Cicero's
education --- the running theme of Q. fr. 3 ---: the rhetor
Paeonius is good, the boy is doing well, but the genre is
the declamatory rather than the theoretical, and Cicero will
draw him into ``our paths'' on country trips. The closing
charge to Quintus to write where and with what hope he is
wintering, the small private hinge that the next letter
(Q. fr. 3.4) will not yet have an answer to.}
